% SOURCE_STATEMENT: "For gradient fields and Gibbs models with non-convex interactions between microscopic particles, what can be proved about the free energy and the Gibbs states?"
% Preparation and final self-audit: CONTRIBUTING.md prompts C and D; source/audit artifacts retained in the local physics_statements_2026_10_11 batch.
% !TeX program = lualatex
\documentclass[11pt]{article}
\usepackage[margin=25mm]{geometry}
\usepackage{fontspec}
\setmainfont{Latin Modern Roman}
\usepackage{amsmath,amssymb,mathtools}
\usepackage{unicode-math}
\setmathfont{Latin Modern Math}
\usepackage{xurl,hyperref}
\hypersetup{colorlinks=true,urlcolor=blue}
\setcounter{secnumdepth}{0}
\setlength{\parindent}{0pt}
\setlength{\parskip}{5pt}
\setlength{\emergencystretch}{3em}
\newcommand{\sref}[2]{\hyperlink{source:#1:#2}{[S#2]}}
\newcommand{\cataloguescope}[1]{\paragraph{Recorded target.}#1}
\newcommand{\R}{\mathbb R}
\newcommand{\C}{\mathbb C}
\newcommand{\N}{\mathbb N}
\newcommand{\Z}{\mathbb Z}
\newcommand{\Q}{\mathbb Q}
\newcommand{\F}{\mathbb F}
\newcommand{\eps}{\varepsilon}
\begin{document}
\section{Nonconvex Gradient Gibbs States and Free Energy}\label{statement}

\subsection{Definitions and mathematical statement}

\paragraph{Shared notation.}$\mathbb N=\{1,2,\ldots\}$; the standard number fields have their usual meanings. All additional conventions are stated locally.


\paragraph{Definitions and conventions.} A height field is $\phi:\mathbb Z^d\to\mathbb R$ and its gradient variables are $\eta_i(x)=\phi(x+e_i)-\phi(x)$. A gradient Gibbs state is a probability measure on these variables satisfying the local Gibbs conditional distributions. Boundary tilt $u\in\mathbb R^d$ controls the mean gradient. The free energy, or surface tension, is the appropriately normalized limit of $-\beta^{-1}\log Z_\Lambda(u)$. Nonconvex interactions have potentials whose Hessian need not be positive, so standard uniformly convex uniqueness and fluctuation arguments need not apply.

\paragraph{Source scope and qualifications.} The report discusses perturbations of a Gaussian potential and a low-temperature RG breakthrough. The question asks what can be proved more broadly; it does not assert universal uniqueness or strict convexity for every nonconvex interaction.

\paragraph{Catalogue mathematical target.}
\paragraph{Catalogue wording retained for traceability.}
For gradient fields and Gibbs models with non-convex interactions between microscopic particles, what can be proved about the free energy and the Gibbs states?

\paragraph{Saved literature context.} The supplied catalogue's assessment dated 2026-08-21 records: Low-temperature renormalization-group methods give first results for free energy and Gibbs states with non-convex interactions, while the general non-convex theory requested in the report remains open. The assessment is retained as dated provenance; it is not a new comprehensive literature or open-status certification.
\subsection{Short English statement}
Develop a general theory of free energy and Gibbs states for nonconvex gradient interactions.
\subsection{Sources}
\begin{itemize}
\item[\hypertarget{source:statement:1}{S1}] ULAM.AI and the UnsolvedMath Contributors, UnsolvedMath, catalogue entry. Catalogue curation under CC BY 4.0; definitions and English exposition added locally.
\url{https://huggingface.co/datasets/ulamai/UnsolvedMath}.
\item[\hypertarget{source:statement:2}{S2}] Interplay of Analysis and Probability in Physics (2008), Oberwolfach Reports, DOI 10.4171/owr/2008/55, pp. 3065--3138; official report source and mathematical context.
\url{https://doi.org/10.4171/owr/2008/55}.
\item[\hypertarget{source:statement:3}{S3}] Official EMS Press PDF of the same Oberwolfach report; PDF page 44.
\url{https://ems.press/content/serial-article-files/46198}.
\item[\hypertarget{source:statement:4}{S4}] non-convex gradient models; reference in the saved dated catalogue literature review.
\url{https://www.mis.mpg.de/events/event/strict-convexity-of-the-surface-tension-of-gradient-models-with-non-convex-interactions}.
\end{itemize}
\label{end:statement}
\end{document}
